\section{DNF Expressions}

A \emph{disjunctive normal form} (DNF) expression is
any sum $m_1 + m_2 + \ldots + m_r$ of monomials where each
monomial $m_i$ is a product of literals. For example
$p_1 p_2 + \bar{p}_1 p_3 p_4$ is a DNF expression. Such expressions appear particularly easy for humans to
comprehend. Hence we expect that any practical
learning system would have to allow for them.

An expression is \emph{monotone} if no variable is
negated in it. We shall show that for monotone DNF
expressions there exists a simple deduction procedure that uses both EXAMPLES and ORACLE. For
unrestricted DNF a similar result can be proved
with respect to a different size measure. In the
unrestricted case there is the additional difficulty
that we can guarantee to deduce a program only for
the function, and not for the concept. This difficulty does not arise in the monotone case where
given an expression we can always compute the value
of the associated concept for a vector $v$ by
making the substitution and asking whether the resulting expression is identically true.

A monomial $m$ is a \emph{prime implicant} of a function $F$ (or of an expression representing the
function) if $m \Rightarrow F$ and if $m' \not\Rightarrow F$ for any $m'$
obtained by deleting one literal from $m$. A DNF
expression is \emph{prime} if it consists of the sum of
prime implicants none of which is redundant in the
sense of being implied by the sum of the others.
There is always a unique prime DNF expression in
the monotone case, but not in general. We therefore define the \emph{degree} of a DNF expression to be
the largest number of monomials that a prime DNF
expression equivalent to it can have. The unique
prime DNF expression in the monotone case is simply
the sum of all the prime implicants.

\medskip
\noindent\textbf{THEOREM B:} \emph{The class of monotone DNF expressions
is learnable via an algorithm B that uses
$L = L(h,d)$ calls of EXAMPLES and \; $dt$ calls of
ORACLE, where $d$ is the degree of the DNF expression $f$ to be learnt and $t$ the number of
variables.}

\medskip
\noindent\textit{Proof.}\quad The algorithm is initialized with formula
$g$ identically equal to the constant zero. The
algorithm then calls EXAMPLES $L$ times to produce
positive examples of $f$. Each time a vector $v$
is produced such that $v \neq g$ a new monomial $m$
is added to $g$. The monomial $m$ is the product
of those literals determined in $v$ that are
essential to make $v \Rightarrow f$. More precisely, the loop
that is executed $L$ times is the following:

\begin{verbatim}
begin v := EXAMPLES
      if v != g then
      begin for i := 1 to t do
                if p_i is determined in v then
                begin set v equal to v but with p_i := *;
                       if ORACLE(v) = 1 then v := v
                end
            set m equal to the product of all
                    literals q such that v => q;
            g := g + m
      end
end
\end{verbatim}

The test $v \neq g$ amounts to asking whether
none of the monomials of $g$ is made true by the
values determined to be true in $v$. Every time
EXAMPLES produces a value $v$ such that $v \neq g$ the
inner loop of the algorithm will find a prime implicant $m$ to add to $g$. Each such $m$ is
different from any previously added (since the
contrary would have implied that $v \Rightarrow g$). It
follows that such a $v$ will be found at most $d$
times and hence ORACLE will be called at most $dt$
times.

Let $X = \Sigma D(w)$ with summation over all (not
necessarily total) vectors $w$ such that $w \Rightarrow f$
but $w \neq g$. This quantity is defined for each
intermediate value of $g$ in the course of the
algorithm, is initially unity, and decreases monotonically with time. Now a monomial will be added
to $g$ each time EXAMPLES outputs a vector $v$ such
that $v \neq g$. The probability of this occurring at
any call of EXAMPLES is exactly $X$. The process
of running the algorithm to completion can have
just two possible outcomes: (i) At some time $X$
has become less than $h^{-1}$, in which case the final
expression $g$ found will approximate to $f$ as
required by the definition of learnability.
(ii) The value of $X$ never goes below $h^{-1}$ and
hence $g$ is not an acceptable approximation. The
probability of the latter eventuality is, however,
at most $h^{-1}$ since it corresponds to the situation
of performing $L(h, (2t)^{k+1})$ Bernoulli experiments
(i.e.\ calls of EXAMPLES) each with probability
greater than $h^{-1}$ of success (i.e.\ of finding a $v$
such that $v \neq g$) and obtaining fewer than $d$
successes (each manifested by the addition of a
new monomial). \hfill $\square$

\medskip

For unrestricted DNF expressions several
problems arise. The main source of difficulty is
the fact that the problem of determining whether
a nontotal vector implies the function specified by
a DNF formula is NP-hard. This is simply
because the problem of determining whether the nowhere determined vector implies the function is the
tautology question of Cook~[3]. An immediate
consequence of this is that it is unreasonable to
assume that the program being learnt computes concepts rather than functions. Another consequence
is that in any algorithm akin to Algorithm B the
test $v \neq g$ may not be feasible if $v$ is not
total. Algorithm B is sufficient, however, to
establish the following:

\medskip
\noindent\textbf{THEOREM B\,$'$:} \emph{Suppose the notion of learnability is
restricted to distributions $D$ such that $D(v) = 0$
whenever $v$ is not total. Then the class of DNF
expressions is learnable, in the sense that programs
for the functions (but not necessarily the concepts)
can be deduced in $L(h,d)$ calls of EXAMPLES and \; $dt$
calls of ORACLE where $d$ is the degree of the DNF
expression to be learnt.}

\medskip

Finally, it may be worth emphasizing that the
monotone case is nonproblematic and the deduction
procedure for it very simple-minded. It may be
interesting to pursue more sophisticated deduction
procedures for it if contexts can be found in
which they can be proved advantageous. The
question as to whether monotone DNF expressions
can be learnt from EXAMPLES alone is open. A positive answer would be especially significant.
